\section{Morse 与 Fukaya \texorpdfstring{$A_\infty$}{A∞} 范畴}\label{sec:8.3}

\SourceParagraph{ch08-8.3-intro-001}{605}
现在我们将逐步引出 Fukaya 范畴的定义。它的特点与此前考虑的范畴有着根本差别。作为准备，我们先给出一个较简单的例子，即 Morse 范畴。它与 Fukaya 范畴具有相似的特点，并且在某些情形下与后者存在联系。

\subsection{Morse 范畴}\label{sec:8.3.1}

\SourceParagraph{ch08-8.3.1-001}{605}
我们先回顾 Morse 理论的基本内容（例如，可参见 \BookBibText{ref242}{MS1} 第 3.4 节和第 10.5 节）。本小节中，$M$ 始终表示一个 $n$ 维紧致定向 Riemann 流形。若可微函数 $f\colon M\to\mathbb R$ 的临界点 $p$ 附近存在局部坐标，使得
\[
f=f(p)+x_1^2+\cdots+x_k^2-x_{k+1}^2-\cdots-x_n^2,
\]
则称该临界点\emph{非退化}。此时，我们称它的指标为 $k$，并记
\[
\deg p=k.
\]

\SourceParagraph{ch08-8.3.1-002}{605}
对于这样的临界点，有两个流形 $W^s(p)$ 和 $W^u(p)$。其中，$W^s(p)$ 为稳定流形，由在 $f$ 的\emph{负}梯度流下、当 $t\to+\infty$ 时趋于 $p$ 的点组成；$W^u(p)$ 为不稳定流形，由在 $f$ 的\emph{正}梯度流下、当 $t\to+\infty$ 时趋于 $p$ 的点组成。由临界点的局部描述，很容易验证
\[
\deg p=\dim W^s(p).
\]

\BookStatementNumber{9}
\begin{definition}\label{def:8.9}
\SourceParagraph{ch08-8.3.1-003}{605}
若 $f$ 只有有限个临界点，这些临界点均非退化，并且所有稳定子流形与不稳定子流形都彼此横截相交，则称 $f$ 为 \emph{Morse--Smale 函数}。
\end{definition}

\SourceParagraph{ch08-8.3.1-004}{606}
对于任意临界点 $p$，$W^s(p)$ 并不自带一个内在的定向。不过，我们可以考虑由 $W^s(p)$ 的两个可能定向生成的自由 Abel 群，再对关系
\[
[\Omega_p]+[-\Omega_p]=0
\]
取商，所得的群记为 $o_p$。这里，$[\Omega_p]$ 和 $[-\Omega_p]$ 表示 $W^s(p)$ 的两个相反定向。作为抽象群，有 $o_p\cong\mathbb Z$。

\SourceParagraph{ch08-8.3.1-005}{606}
现在可以定义 Morse 复形
\[
CM^i(f)=\bigoplus_{\substack{p\in\operatorname{Crit}(f)\\\deg(p)=i}}o_p.
\]
我们按照如下方式定义微分，并将它记为 $m_1$。实质上，$m_1$ 计数连接两个次数相差 $1$ 的临界点的梯度流线。具体来说，首先注意，$M$ 已定向，且 $W^s(p)$ 与 $W^u(p)$ 在 $p$ 处横截相交，维数互补，因此 $W^s(p)$ 的一个定向会诱导 $W^u(p)$ 的定向。若 $\deg q=\deg p+1$，则
\[
\dim W^u(p)=n-\deg p,\qquad \dim W^s(q)=\deg p+1,
\]
所以 $W^u(p)\cap W^s(q)$ 为一维，因而是一些线段的并：这些线段就是从 $p$ 到 $q$ 的 $-f$ 的梯度流轨道。此外，$W^u(p)$ 和 $W^s(q)$ 的定向决定了 $W^u(p)\cap W^s(q)$ 的定向。映射 $f$ 将这些线段与 $\mathbb R$ 中的线段对应起来，而 $\mathbb R$ 具有自然定向。我们令 $n(q,p)$ 等于 $W^u(p)\cap W^s(q)$ 中定向与 $\mathbb R$ 的定向一致的线段数，减去定向相反的线段数。由此定义
\[
m_1\colon CM^i(f)\longrightarrow CM^{i+1}(f),
\qquad
m_1([\Omega_p])=\sum_q n(q,p)[\Omega_q].
\]
注意，改变 $\Omega_p$ 或 $\Omega_q$ 的定向，会改变 $W^u(p)\cap W^s(q)$ 的定向符号，因而改变 $n(q,p)$ 的符号。因此 $m_1$ 定义良好。

\SourceParagraph{ch08-8.3.1-006}{606}
\BookBibText{ref242}{MS1} 第 10.5.4 节已经证明，映射 $m_1$ 满足 $m_1^2=0$。我们在这里的目标是定义具有高阶乘法映射的 $A_\infty$ 范畴，而 $m_1^2=0$ 将被包含在 $A_\infty$ 关系之中。

\SourceParagraph{ch08-8.3.1-007}{606}
现在来定义 \emph{Morse 范畴} $\operatorname{Morse}(M)$，其对象为 $M$ 上的所有光滑实值函数。Fukaya 在\BookCite{ref165}中描述过这一构造，不过这里遵循 Abouzaid 在\BookCite{ref002}中的阐述。

\SourceParagraph{ch08-8.3.1-008}{606}
$\operatorname{Morse}(M)$ 的对象为 $M$ 上的所有光滑实值函数。但是，态射空间和高阶乘法映射并非总是存在。因此，我们所定义的甚至还不是一个 $A_\infty$ 范畴，而是一个称为 $A_\infty$ \emph{预范畴}的结构；下一小节将更精确地说明它的含义。

\SourceParagraph{ch08-8.3.1-009}{607}
若两个函数 $f_0,f_1$ 满足 $f_1-f_0$ 是 $M$ 上的 Morse--Smale 函数，则令
\[
\operatorname{Hom}(f_0,f_1)=CM^\bullet(f_1-f_0),
\]
其中 $m_1$ 如上定义。给定满足某些附加条件的 $\{f_0,\ldots,f_d\}$——这些条件稍后会准确说明——我们将定义高阶乘法映射
\[
m_d\colon\operatorname{Hom}(f_{d-1},f_d)\otimes\cdots\otimes
\operatorname{Hom}(f_0,f_1)\longrightarrow\operatorname{Hom}(f_0,f_d).
\]
为此，进行如下构造。

\BookStatementNumber{10}
\begin{definition}\label{def:8.10}
\SourceParagraph{ch08-8.3.1-010}{607}
给定一列函数 $f_0,\ldots,f_d$，以及 $f_{i+1}-f_i$ 的临界点 $p_{i,i+1}$ 和 $f_d-f_0$ 的临界点 $p_{0,d}$，\emph{梯度树模空间}
\[
\mathcal S_d(p_{0,d};p_{0,1},\ldots,p_{d-1,d})
\]
定义为满足以下条件的映射 $\phi\colon S\to M$ 所组成的空间，其中 $S\in\mathcal S_d$，后者为带状树模空间。
\begin{enumerate}
\item 按如下方式为 $S$ 的每条边 $e$ 标记一个函数 $f_e$。若 $e$ 是标号为 $1\leq i\leq d$ 的外部输入边，则 $f_e=f_i-f_{i-1}$；否则，若 $e$ 从顶点 $v$ 输出，则 $f_e$ 等于所有输入 $v$ 的边所标记的函数之和。每条边 $e$ 的像均须为 $f_e$ 的一条梯度流线。
\item $e$ 的定向由 $-f_e$ 的梯度流给出。此外，若 $e$ 的长度为 $l$，则 $\phi$ 将由边的一个固定参数化 $[0,l]$ 给出的向量场，与 $-f_e$ 的梯度向量场对应起来。
\item 若 $e$ 是外部边，且 $f_e=f_i-f_j$，则 $\phi$ 的像在外部顶点处趋于 $p_{j,i}$。
\end{enumerate}
\end{definition}

\SourceParagraph{ch08-8.3.1-011}{607}
对于给定的函数列 $f_0,\ldots,f_d$，可以令
\[
\mathcal S(f_0,\ldots,f_d)=
\coprod_{p_{0,d},p_{0,1},\ldots,p_{d-1,d}}
\mathcal S(p_{0,d};p_{0,1},\ldots,p_{d-1,d}),
\]
这里的不交并遍历所有可能的临界点。

\SourceParagraph{ch08-8.3.1-012}{607}
允许边的长度趋于 $\infty$，便可将此模空间紧化；此时，一棵梯度树趋于若干梯度树的并，如\BookRef{fig:8.16}{图16}所示。这样得到紧化 $\overline{\mathcal S}(f_0,\ldots,f_d)$。

\setcounter{figure}{15}
\begin{figure}[H]
\centering
\BookImage{01_Numbered_Figures/Chapter_08/Ch08_Fig16_PDFp608_Bookp598.pdf}
\BookNumberOnlyCaption\label{fig:8.16}
\end{figure}

\SourceParagraph{ch08-8.3.1-013}{607}
下面说明如何构造 $\mathcal S(p_{0,d};p_{0,1},\ldots,p_{d-1,d})$。记 $p_i:=p_{i-1,i}$、$q=p_{0,d}$，并令
\[
\mathcal E(q;p_1,\ldots,p_d)=W^s_{f_d-f_0}(q)
\times\left(\prod_{i=1}^dW^u_{f_i-f_{i-1}}(p_i)\right)\times\mathcal S_d,
\]
其中 $W$ 的下标表示我们针对哪个函数取稳定流形或不稳定流形。对每个 $S\in\mathcal S_d$，在输出边上选取一个基点 $s\in S$，使它与最后一个内部顶点相距某个固定距离。于是，对于给定的 $S$，可以按如下方式定义映射
\[
\Phi_j^s\colon W^u_{f_j-f_{j-1}}(p_j)\longrightarrow M.
\]

\SourceParagraph{ch08-8.3.1-014}{608}
令 $v_j$ 为第 $j$ 条输入边所接到的内部顶点，$A$ 为 $S$ 中从 $v_j$ 到 $s$ 的路径。对于点 $y\in W^u_{f_j-f_{j-1}}(p_j)$，依次施加 $A$ 所经过的各条边上函数的梯度流，取 $y$ 在这组梯度流下的像；其中，$-f_e$ 的梯度流所经历的时间等于 $e\cap A$ 的长度。将这些映射与包含映射 $W^s_{f_d-f_0}(q)\hookrightarrow M$ 合在一起，得到具有 $d+1$ 个分量的映射
\[
\Phi\colon\mathcal E(q;p_1,\ldots,p_d)\longrightarrow M^{d+1}.
\]

\SourceParagraph{ch08-8.3.1-015}{608}
此时很明显，$\mathcal S(q;p_1,\ldots,p_d)$ 正是对角线 $\Delta\subseteq M^{d+1}$ 的原像。若每个 $f_j-f_i$ 在 $i<j$ 时均为 Morse--Smale 函数，且 $\Phi$ 与 $\Delta$ 横截，即 $\Phi^{-1}(\Delta)$ 是具有预期维数的光滑流形，则称函数组 $f_0,\ldots,f_d$ 为 \emph{Morse--Smale 函数组}。注意，这一预期维数为
\[
\begin{aligned}
&\dim\mathcal E(q;p_1,\ldots,p_d)-\operatorname{codim}_{M^{d+1}}\Delta\\
&\quad=\deg q+d\dim M-\sum_i\deg p_i+d-2-d\dim M\\
&\quad=\deg q+d-2-\sum_i\deg p_i.
\end{aligned}
\]
因此，特别地，当 $2-d+\sum_i\deg p_i=\deg q$ 时，$\overline{\mathcal S}(f_0,\ldots,f_d)$ 是零维的。

\SourceParagraph{ch08-8.3.1-016}{608}
这里有一个微妙之处：我们需要为 $\overline{\mathcal S}(f_0,\ldots,f_d)$ 赋予定向。这一步必须格外谨慎，才能使 $A_\infty$ 关系在正确的符号约定下成立。只有\BookCite{ref002}附录 B 详细完成了这项工作。这里略去这些细节。

\SourceParagraph{ch08-8.3.1-017}{609}
当 $2-d+\sum_i\deg p_i=\deg q$ 时，令 $n(q;p_1,\ldots,p_d)$ 为 $\mathcal S(q;p_1,\ldots,p_d)$ 中梯度树的带符号数目，然后用公式
\[
m_d([\Omega_{p_d}]\otimes\cdots\otimes[\Omega_{p_1}])
=\sum_q n(q;p_1,\ldots,p_d)[\Omega_q]
\]
定义 $m_d$。当然，定向 $\Omega_{p_1},\ldots,\Omega_{p_d},\Omega_q$ 会影响模空间 $\mathcal S(q;p_1,\ldots,p_d)$ 的定向，从而使此定义良好。

\SourceParagraph{ch08-8.3.1-018}{609}
还需要说明，如此定义的 $m_d$ 为什么满足 $A_\infty$ 关系。为此，考虑满足
\[
3-d+\sum_{i=1}^d\deg p_i=\deg q
\]
时的紧化 $\overline{\mathcal S}(q;p_1,\ldots,p_d)$。由这一次数条件可知，模空间是一维的。其边界由形如
\[
\overline{\mathcal S}(q;p_1,\ldots,p_j,r,p_{i+j+1},\ldots,p_d)
\times\overline{\mathcal S}(r;p_{j+1},\ldots,p_{i+j})
\]
的部分组成，其中 $1\leq i\leq d$、$0\leq j\leq d-i$，并且
\[
\deg r=2-i+\sum_{k=j+1}^{i+j}\deg p_k.
\]

\SourceParagraph{ch08-8.3.1-019}{609}
现在考虑 $d$ 次 $A_\infty$ 关系中出现的一项：
\[
\begin{aligned}
&m_{d-i+1}\bigl([\Omega_{p_d}],\ldots,[\Omega_{p_{i+j+1}}],
m_i([\Omega_{p_{i+j}}],\ldots,[\Omega_{p_{j+1}}]),
[\Omega_{p_j}],\ldots,[\Omega_{p_1}]\bigr)\\
&\quad=\sum_r n(r;p_{j+1},\ldots,p_{i+j})\,
m_{d-i+1}\bigl([\Omega_{p_d}],\ldots,[\Omega_{p_{i+j+1}}],
[\Omega_r],[\Omega_{p_j}],\ldots,[\Omega_{p_1}]\bigr)\\
&\quad=\sum_{q,r}n(r;p_{j+1},\ldots,p_{i+j})\,
n(q;p_1,\ldots,p_j,r,p_{i+j+1},\ldots,p_d)[\Omega_q].
\end{aligned}
\]

\SourceParagraph{ch08-8.3.1-020}{609}
因此，临界点 $p_1,\ldots,p_d$ 对 $A_\infty$ 关系中 $[\Omega_q]$ 的系数所作的贡献，恰好是 $\overline{\mathcal S}(q;p_1,\ldots,p_d)$ 的边界点数的带符号和。然而，一维带边界流形的边界点数的带符号和总是零，所以只要所有定向选择彼此一致，$A_\infty$ 方程就成立。\BookCite{ref002}描述了这种定向的选择。

\subsection{\texorpdfstring{$A_\infty$}{A∞} 预范畴}\label{sec:8.3.2}

\SourceParagraph{ch08-8.3.2-001}{609}
上一小节的构造存在一个缺点：并非任意两个对象之间都定义了态射空间，而且高阶乘法还受到进一步的限制。因此，我们甚至尚未定义一个真正的 $A_\infty$ 范畴，而只是得到了 Kontsevich 与 Soibelman 在\BookCite{ref311}中称为 $A_\infty$ \emph{预范畴}的结构。下面给出定义。

\BookStatementNumber{11}
\begin{definition}\label{def:8.11}
\SourceParagraph{ch08-8.3.2-002}{609}
一个（非幺）$A_\infty$ 预范畴 $\mathcal A$ 包含以下数据：
\begin{enumerate}
\item 一个对象类 $\operatorname{Ob}(\mathcal A)$。
\item 对每个 $n\geq2$，给定一族\emph{横截序列}
$\operatorname{Ob}_{\mathrm{tr}}^n(\mathcal A)\subseteq
\operatorname{Ob}(\mathcal A)^n$，即由对象的 $n$ 元序列组成的一个集合。
\item 对每个 $(X_0,X_1)\in\operatorname{Ob}_{\mathrm{tr}}^2(\mathcal A)$，给定一个 $\mathbb Z$ 分次链复形 $\operatorname{Hom}_{\mathcal A}(X_0,X_1)$。
\item 对每个 $(X_0,\ldots,X_d)\in\operatorname{Ob}_{\mathrm{tr}}^{d+1}(\mathcal A)$，给定映射
\[
\begin{aligned}
m_d\colon\operatorname{Hom}_{\mathcal A}(X_{d-1},X_d)\otimes\cdots
\otimes\operatorname{Hom}_{\mathcal A}(X_0,X_1)
\longrightarrow\operatorname{Hom}_{\mathcal A}(X_0,X_d)[2-d].
\end{aligned}
\]
此外，还要求：
\item 横截序列的每个子序列都是横截的。
\item 各映射 $m_d$ 满足 $A_\infty$ 关系。
\end{enumerate}
\end{definition}

\SourceParagraph{ch08-8.3.2-003}{610}
因此，$\operatorname{Morse}(M)$ 是一个 $A_\infty$ 预范畴的例子：它的横截序列为那些自身以及每个子序列都满足 Morse--Smale 条件的序列 $(f_0,\ldots,f_d)$。

\SourceParagraph{ch08-8.3.2-004}{610}
我们还需要 $A_\infty$ 预范畴之间的拟等价概念。首先定义 $A_\infty$ 预范畴之间的函子。

\BookStatementNumber{12}
\begin{definition}\label{def:8.12}
\SourceParagraph{ch08-8.3.2-005}{610}
设 $\mathcal A,\mathcal B$ 为 $A_\infty$ 预范畴。一个 $A_\infty$ 函子 $F\colon\mathcal A\to\mathcal B$ 包含：
\begin{enumerate}
\item 映射 $F\colon\operatorname{Ob}(\mathcal A)\to\operatorname{Ob}(\mathcal B)$，满足只要 $(X_1,\ldots,X_d)\in\operatorname{Ob}_{\mathrm{tr}}^d(\mathcal A)$，就有
\[
(F(X_1),\ldots,F(X_d))\in\operatorname{Ob}_{\mathrm{tr}}^d(\mathcal B).
\]
\item 对 $\mathcal A$ 的每个横截序列 $(X_0,\ldots,X_d)$，给定次数为 $1-d$ 的映射
\[
\begin{aligned}
F_d\colon\operatorname{Hom}_{\mathcal A}(X_{d-1},X_d)\otimes\cdots
\otimes\operatorname{Hom}_{\mathcal A}(X_0,X_1)
\longrightarrow\operatorname{Hom}_{\mathcal B}(F(X_0),F(X_d)),
\end{aligned}
\]
这些映射满足 $A_\infty$ 函子的通常公式。
\end{enumerate}
\end{definition}

\BookStatementNumber{13}
\begin{definition}\label{def:8.13}
\SourceParagraph{ch08-8.3.2-006}{610}
若 $A_\infty$ 预范畴 $\mathcal A,\mathcal B$ 之间的 $A_\infty$ 函子在 $H^*(\mathcal A)$、$H^*(\mathcal B)$ 层次诱导同构，则称它为\emph{拟等价}。
\end{definition}

\SourceParagraph{ch08-8.3.2-007}{610}
在许多情形下，例如即将在\BookSectionRef{sec:8.4}{8.4}节看到的情形，避开对象之间不横截时的定义通常比较方便；这一形式体系使我们能够这样处理。

\subsection{Fukaya 范畴}\label{sec:8.3.3}

\SourceParagraph{ch08-8.3.3-001}{610}
现在描述 Calabi--Yau 流形 $X$ 的 Fukaya 范畴，其中 $X$ 配备一个辛形式 $\omega$ 和一个处处非零的全纯 $n$ 形式 $\Omega$。我们将比本书前面的介绍（见\BookSectionRef{sec:3.6}{3.6}节）以及 \BookBibText{ref242}{MS1} 中的介绍稍微详细一些。不过，要尝试给出在技术上完全正确的定义，远远超出了这里所能完成的范围。有关更多技术细节，可查阅 Seidel 的书\BookCite{ref421}；这里紧密遵循其中的处理。

\SourceParagraph{ch08-8.3.3-002}{610}
首先描述 Fukaya 范畴 $\operatorname{Fuk}(X)$ 的对象。它们是四元组 $(L,\mathcal L,\widetilde\xi_L,S_L)$，其中：
\begin{itemize}
\item $L\subseteq X$ 是一个定向 Lagrange 子流形。
\item $\mathcal L$ 是 $L$ 上具有平坦联络的 $U(1)$ 丛。
\item 总存在一个函数
\[
\xi_L\colon L\longrightarrow S^1,\qquad
\xi_L=\frac{\Omega|_L}{|\Omega|_L|}.
\]
换言之，在任意点 $x\in L$ 处，选取 $L$ 的切向量的一组正定向基 $\eta_1,\ldots,\eta_n$，则
\[
\xi_L(x)=\frac{\Omega(\eta_1,\ldots,\eta_n)}
{|\Omega(\eta_1,\ldots,\eta_n)|}.
\]
很容易验证，$\xi_L$ 与基的选择无关。$\widetilde\xi_L$ 是 $L$ 的一个\emph{分次}，即一个提升
\[
\widetilde\xi_L\colon L\longrightarrow\mathbb R,
\qquad \xi_L=e^{\pi\mathrm i\widetilde\xi_L}.
\]

\SourceParagraph{ch08-8.3.3-003}{611}
注意，并非所有 Lagrange 子流形都具有分次。函数 $\xi_L$ 决定了 $H^1(L,\mathbb Z)$ 中的一个类，即沿 $\xi_L$ 拉回 $H^1(S^1,\mathbb Z)$ 的生成元所得的类，称为 $L$ 的 \emph{Maslov 类}，见\BookRef{fig:8.17}{图17}。只有当这个类为零时，分次才存在；分次并不唯一，可以加上任意偶整数。注意，改变 $L$ 的定向需要将分次平移 $1$，这可以理解为把膜变为反膜。若 $L$ 实际上是特殊 Lagrange 子流形，即对于某个固定的角 $\theta$，有 $\Omega|_L=e^{\mathrm i\theta}\operatorname{Vol}(L)$，则 $\xi_L$ 是常函数 $e^{\mathrm i\theta}$。因此，特殊 Lagrange 子流形总是具有分次。

\SourceParagraph{ch08-8.3.3-004}{611}
\item $S_L$ 是 $L$ 上的一个自旋结构。$L$ 上存在自旋结构的障碍是 $L$ 的第二 Stiefel--Whitney 类 $w_2\in H^2(L,\mathbb Z/2\mathbb Z)$；若 $w_2=0$，则自旋结构由 $H^1(L,\mathbb Z/2\mathbb Z)$ 分类。这一选择在我们的阐述中不会起很大的作用，但对于确定那些用来定义映射 $m_k$ 的 $J$ 全纯圆盘模空间的定向，它实际上很重要。由于这里将略去对这些符号的大部分讨论，读者不必过于担心这一问题。
\end{itemize}

\SourceParagraph{ch08-8.3.3-005}{611}
事实上，还需要对 Lagrange 子流形施加进一步的限制：$L$ 应当在\BookCite{ref164}的意义下\emph{无障碍}。这是一个极其精细而复杂的条件，涉及边界位于 $L$ 内的全纯圆盘。稍后会再作一点简短讨论，因为这一问题有大量技术细节。这里只需指出，在我们考虑的情形中，它仅在 $\dim X\geq3$ 时成为问题。

\SourceParagraph{ch08-8.3.3-006}{611}
假设给定 $\operatorname{Fuk}(X)$ 的两个对象，为方便起见，将它们简记为 $L_0$、$L_1$。假设 $L_0$ 与 $L_1$ 横截相交。我们希望定义 $\operatorname{Hom}_{\operatorname{Fuk}(X)}(L_0,L_1)$。即使仅将它定义为一个分次向量空间，也需要几个概念，更不用说把它定义为一个复形了。

\begin{figure}[H]
\centering
\BookImage{01_Numbered_Figures/Chapter_08/Ch08_Fig17_PDFp612_Bookp602.pdf}
\caption{曲面 $X$ 内的一条可缩闭环 $L$ 总具有非平凡的 Maslov 类：在图示例子中，切线的角度从 $0$ 变化到 $2\pi$。}\label{fig:8.17}
\end{figure}

\SourceParagraph{ch08-8.3.3-007}{611}
首先，需要给出交点 $x\in L_0\cap L_1$ 的指标的正确概念。具体如下。考虑辛向量空间 $V:=T_xX$，以 $\operatorname{Gr}(V)$ 表示 Lagrange Grassmann 流形，即 $V$ 的所有 Lagrange 子空间组成的空间。因此，对于 $i=0,1$，$\Lambda_i:=T_xL_i$ 给出了 $\operatorname{Gr}(V)$ 中的两个点。由横截性可知，$\Lambda_0\cap\Lambda_1=0$。若 $\operatorname{Gr}(V)$ 中连接 $\Lambda_0$ 与 $\Lambda_1$ 的路径 $\Lambda_t$（$0\leq t\leq1$）满足 $\Lambda_0\cap\Lambda_t=0$，$0<t\leq1$，且在 $t=0$ 处的\emph{穿越形式}负定，则称这条路径为\emph{无穿越路径}。这一 $\Lambda_0$ 上的二次型如下定义：对充分小的 $t$，先选择同构 $\phi_t\colon\Lambda_0\to\Lambda_t$，使 $\phi_0$ 为恒等映射，再令二次型为
\[
v\longmapsto\left.\frac{d}{dt}\omega(v,\phi_t(v))\right|_{t=0}.
\]

\SourceParagraph{ch08-8.3.3-008}{612}
例如，总可以取 $V=\mathbb C^n$，配备标准辛形式，令 $\Lambda_0=\mathbb R^n$，并假设
\[
\Lambda_1=e^{\mathrm i\pi c_1}\mathbb R\times\cdots\times
e^{\mathrm i\pi c_n}\mathbb R,\qquad c_k\in(-1,0).
\]
那么
\[
\Lambda_t=e^{\mathrm i\pi t c_1}\mathbb R\times\cdots\times
e^{\mathrm i\pi t c_n}\mathbb R
\]
就是连接 $\Lambda_0$ 与 $\Lambda_1$ 的一条无穿越路径。此时，总可以选择函数 $\widetilde a\colon[0,1]\to\mathbb R$，使
\[
e^{2\pi\mathrm i\widetilde a(t)}
=\left(\frac{\Omega|_{\Lambda_t}}{|\Omega|_{\Lambda_t}|}\right)^2.
\]
然后，将交点 $x\in L_0\cap L_1$ 的\emph{绝对 Maslov 指标}定义为
\[
I(x):=(\widetilde\xi_{L_1}(x)-\widetilde a(1))
-(\widetilde\xi_{L_0}(x)-\widetilde a(0)).
\]
这个数实际上总属于 $\mathbb Z$。\BookRef{fig:8.18}{图18}给出了一个例子。

\begin{figure}[H]
\centering
\BookImage{01_Numbered_Figures/Chapter_08/Ch08_Fig18_PDFp613_Bookp603.pdf}
\caption{$\mathbb C$ 中的两个例子。箭头表示定向，分次为常数 $\widetilde\xi\in(-1,1)$。第一种情形的指标为 $1$，第二种情形的指标为 $0$。}\label{fig:8.18}
\end{figure}

\SourceParagraph{ch08-8.3.3-009}{612}
其次，由于收敛性问题，不能只在基域 $\mathbb C$ 上工作，而要使用 \emph{Novikov 环} $\Lambda_{\mathrm{nov}}$。它定义为形如
\[
\sum_{i\in\mathbb Z}a_iq^{\lambda_i}
\]
的形式幂级数所组成的环，其中系数 $a_i\in\mathbb Z$ 对所有充分小的负整数 $i$ 均为零，而实数列 $\lambda_i$ 满足
\[
\lim_{i\to\infty}\lambda_i=\infty.
\]
后一条件保证这个环中的乘法有意义。

\SourceParagraph{ch08-8.3.3-010}{613}
当 $L_0$ 与 $L_1$ 横截相交时，现在可以定义 Floer 复形 $\operatorname{Hom}(L_0,L_1)$。暂且限制于局部系统 $\mathcal L_0,\mathcal L_1$ 平凡的情形，稍后再考虑一般情形。令
\[
\operatorname{Hom}^i(L_0,L_1)=
\bigoplus_{\substack{x\in L_0\cap L_1\\I(x)=i}}[x]\cdot\Lambda_{\mathrm{nov}}.
\]

\SourceParagraph{ch08-8.3.3-011}{613}
利用伪全纯带状域的模空间来定义微分 $m_1$。取一族与辛形式 $\omega$ 相容的殆复结构 $J:=\{J_t\}_{0\leq t\leq1}$，即要求 $\omega(\cdot,J_t\cdot)$ 在 $X$ 上定义 Riemann 度量。对于 $p,q\in L_0\cap L_1$，以 $\mathcal M(p,q,L_0,L_1,J)$ 表示满足下列条件的映射 $u\colon\mathbb R\times[0,1]\to X$ 的模空间：
\[
\begin{gathered}
u(\mathbb R\times\{0\})\subseteq L_0,\qquad
u(\mathbb R\times\{1\})\subseteq L_1,\\
\frac{\partial u}{\partial\tau}+J_t(u)\frac{\partial u}{\partial t}=0,\\
\lim_{\tau\to-\infty}u(\tau,t)=p,\qquad
\lim_{\tau\to+\infty}u(\tau,t)=q.
\end{gathered}
\]
这里，$\tau$ 是 $\mathbb R$ 上的坐标，$t$ 是 $[0,1]$ 上的坐标。第二行应理解为 Cauchy--Riemann 方程。对于一般选取的 $J_t$，预期有
\[
\dim\mathcal M(p,q,L_0,L_1,J)=I(q)-I(p).
\]
这是通过计算某些算子的指标得到的，属于伪全纯曲线理论中的标准计算。这里不提供这些技术细节。

\SourceParagraph{ch08-8.3.3-012}{613}
注意，$\mathbb R$ 通过平移坐标 $\tau$ 作用于此模空间：对于 $r\in\mathbb R$，把解 $u(\cdot,\cdot)$ 替换为解 $u(\cdot+r,\cdot)$。这个 $\mathbb R$ 作用只有在 $p=q$ 且 $u$ 为常值映射时才不是自由的。于是得到商空间
\[
\mathcal M(p,q,L_0,L_1,J)/\mathbb R,
\]
其维数为 $I(q)-I(p)-1$；当 $p=q$ 时，约定该商为空。可以为这一空间赋予定向，这与我们对 $L_i$ 的自旋结构所作的选择有关。因此，当 $I(q)=I(p)+1$ 时，$\mathcal M(p,q,L_0,L_1,J)/\mathbb R$ 是一个有限集合，每个点 $u$ 都带有符号 $(-1)^{s(u)}$。

\SourceParagraph{ch08-8.3.3-013}{614}
与 Morse 范畴的情形一样，我们利用上述模空间定义 $m_1$。具体而言，若 $p,r\in L_0\cap L_1$，且 $I(r)=I(p)+1$，则 $m_1([p])$ 中 $[r]$ 的系数为
\[
\sum_{u\in\mathcal M(p,r,L_0,L_1,J)/\mathbb R}
(-1)^{s(u)}q^{\int u^*\omega}\in\Lambda_{\mathrm{nov}}.
\]

\SourceParagraph{ch08-8.3.3-014}{614}
若要证明 $m_1^2=0$，理想的做法是仿照 Morse 情形，在 $I(q)=I(p)+2$ 时考虑一维模空间 $\mathcal M(p,q,L_0,L_1,J)/\mathbb R$，再通过研究其边界得到 $m_1^2=0$。然而，与 Morse 情形不同，全纯带状域可以以几种不同方式退化：它可以退化为一个从 $p$ 连接到 $r$ 的带状域与一个从 $r$ 连接到 $q$ 的带状域的并，其中 $I(r)=I(p)+1$；也可能有一个全纯圆盘冒泡出来，其边界包含在 $L_0$ 或 $L_1$ 中，见\BookRef{fig:8.19}{图19}。（熟悉伪全纯曲线通常情形的读者，可能还会担心球面冒泡，但那是实余维为二的现象。）第一类退化提供关于 $m_1^2$ 的信息，而第二类退化是新的，可能导致 $m_1^2=0$ 的障碍。

\begin{figure}[H]
\centering
\BookImage{01_Numbered_Figures/Chapter_08/Ch08_Fig19_PDFp614_Bookp604.pdf}
\caption{全纯带状域的两种可能退化。第二种情形的绘图只是示意：圆盘的边界仍位于其中一个 Lagrange 子流形上。}\label{fig:8.19}
\end{figure}

\SourceParagraph{ch08-8.3.3-015}{615}
\BookCite{ref164}对这个问题进行了十分深入的研究。若边界位于 Lagrange 子流形 $L_0$、$L_1$ 上的全纯圆盘模空间满足某些性质——\BookSectionRef{sec:3.6.2.3}{3.6.2.3}节曾提示这些性质——则可以修改 $m_1$ 的定义，使 $m_1^2=0$。若 Lagrange 子流形 $L$ 不满足相应的性质，就称它\emph{有障碍}，不应将其视为 Fukaya 范畴的对象。这是对 Fukaya 范畴对象的一项附加限制。

\SourceParagraph{ch08-8.3.3-016}{615}
一般而言，可能有许多全纯圆盘的边界位于某个 Lagrange 子流形上，这自然引出了开 Gromov--Witten 不变量。目前尚无完全令人满意的开 Gromov--Witten 不变量数学理论，不过在许多有趣的情形下已有相应理论，见\BookCite{ref297}、\BookCite{ref429}。这里不展开这个主题，只简要讨论这种圆盘模空间的预期维数。

\SourceParagraph{ch08-8.3.3-017}{615}
给定全纯圆盘 $u\colon D\to X$，满足 $u(\partial D)\subseteq L$。对于任意辛流形 $X$，都可以定义 $u$ 的 Maslov 类 $\mu(u)$。具体如下：将丛 $u^*(TX)$ 作为辛向量丛平凡化为 $D\times\mathbb C^n$，其中 $\mathbb C^n$ 带有标准辛结构。这就定义了映射 $\gamma\colon\partial D\to\operatorname{Gr}(\mathbb C^n)$：它把 $x\in\partial D$ 映为在此平凡化下的 $T_{u(x)}L\subseteq\mathbb C^n$。与前面一样，$\mathbb C^n$ 上的标准全纯 $n$ 形式 $\Omega$ 定义映射 $\operatorname{Gr}(\mathbb C^n)\to S^1$，将由 $\eta_1,\ldots,\eta_n$ 张成的子空间映为
\[
\left(\frac{\Omega(\eta_1,\ldots,\eta_n)}
{|\Omega(\eta_1,\ldots,\eta_n)|}\right)^2.
\]
$H^1(S^1,\mathbb Z)$ 的正生成元拉回至 $H^1(\operatorname{Gr}(\mathbb C^n),\mathbb Z)$ 所得的类，称为 \emph{Maslov 类} $\mu$；它生成 $H^1(\operatorname{Gr}(\mathbb C^n),\mathbb Z)$。于是，$u$ 的 Maslov 类为
\[
\mu(u):=\int_{\partial D}\gamma^*\mu\in\mathbb Z.
\]

\SourceParagraph{ch08-8.3.3-018}{615}
若 $J$ 是一般的殆复结构，以 $\mathcal M(L,J)$ 表示边界位于 $L$ 内的 $J$ 全纯映射 $u\colon D\to X$ 的空间。在暂不考虑多重覆盖问题时，预期 $\mathcal M(L,J)$ 中包含 $u$ 的分支具有维数 $\tfrac12\dim_{\mathbb R}X+\mu(u)$。在我们的情形中，若 $L$ 具有分次，则实际上 $\mu(u)=0$。由于 $\operatorname{PSL}_2(\mathbb R)$ 通过重新参数化作用于 $\mathcal M(L,J)$，可得
\[
\dim\mathcal M(L,J)/\operatorname{PSL}_2(\mathbb R)
=\tfrac12\dim_{\mathbb R}X-3.
\]
因此，在低维情况下不必担心这种情形；只有当 Calabi--Yau 流形 $X$ 的复维数至少为 $3$ 时，它才成为问题。

\SourceParagraph{ch08-8.3.3-019}{615}
从现在起，假设我们处于不会出现边界位于某个 Lagrange 子流形上的非恒定全纯圆盘的情形。当然，这排除了我们最感兴趣的三维情形，但我们仍将在这个假设下继续。此时便能够定义
\[
m_d\colon\operatorname{Hom}(L_{d-1},L_d)\otimes\cdots\otimes
\operatorname{Hom}(L_0,L_1)\longrightarrow\operatorname{Hom}(L_0,L_d)[2-d].
\]

\SourceParagraph{ch08-8.3.3-020}{616}
对于 $p_{i-1,i}\in L_{i-1}\cap L_i$（$1\leq i\leq d$）、$p_{0,d}\in L_0\cap L_d$，以及某个一般的殆复结构 $J$，定义
\[
\mathcal M(p_{0,d};p_{0,1},\ldots,p_{d-1,d};L_0,\ldots,L_d)
\]
为如下映射 $u\colon D\to X$ 的模空间：$D$ 是一个圆盘，其边界上有 $d+1$ 个标记点 $x_{0,1},\ldots,x_{d-1,d},x_{0,d}\in\partial D$，按逆时针循环顺序排列，并满足
\[
\begin{aligned}
u(x_{i-1,i})&=p_{i-1,i},&&1\leq i\leq d,\\
u(x_{0,d})&=p_{0,d},\\
u([x_{i-1,i},x_{i,i+1}])&\subseteq L_i,&&1\leq i\leq d-1,\\
u([x_{d-1,d},x_{0,d}])&\subseteq L_d,\\
u([x_{0,d},x_{0,1}])&\subseteq L_0.
\end{aligned}
\]
这里，对于 $D$ 上相邻的两个标记点 $x,y$，$[x,y]$ 表示 $\partial D$ 上连接它们的弧段。这是全纯多边形的模空间，见\BookRef{fig:8.20}{图20}。

\begin{figure}[H]
\centering
\BookImage{01_Numbered_Figures/Chapter_08/Ch08_Fig20_PDFp616_Bookp606.pdf}
\BookNumberOnlyCaption\label{fig:8.20}
\end{figure}

\SourceParagraph{ch08-8.3.3-021}{616}
对于给定的带有 $d+1$ 个标记点的圆盘，此模空间的预期维数为 $I(p_{0,d})-\sum_{i=1}^d I(p_{i-1,i})$。由于这类带有 $d+1$ 个标记点的圆盘的模空间为维数 $d-2$ 的 Stasheff 结合多面体 $\mathcal S_d$，得到
\[
\begin{aligned}
&\dim\mathcal M(p_{0,d};p_{0,1},\ldots,p_{d-1,d};L_0,\ldots,L_d)
/\operatorname{PSL}_2(\mathbb R)\\
&\hspace{4em}=I(p_{0,d})+d-2-\sum_{i=1}^d I(p_{i-1,i}).
\end{aligned}
\]
当这个维数为零时，$m_d([p_{d-1,d}],\ldots,[p_{0,1}])$ 中 $[p_{0,d}]$ 的系数为
\[
\sum_{u\in\mathcal M(p_{0,d};p_{0,1},\ldots,p_{d-1,d};L_0,\ldots,L_d)
/\operatorname{PSL}_2(\mathbb R)}
(-1)^{s(u)}q^{\int u^*\omega}\in\Lambda_{\mathrm{nov}}.
\]
这里，$s$ 仍表示那个尚未展开解释的符号选择，它依赖于自旋结构的选择。

\SourceParagraph{ch08-8.3.3-022}{617}
假设全纯圆盘冒泡不会造成问题，则这些 $m_k$ 满足 $A_\infty$ 关系。证明仍然采用通常的方法，即分析一维圆盘模空间的紧化方式。

\SourceParagraph{ch08-8.3.3-023}{617}
现在说明，当 Lagrange 子流形 $L_i$ 带有非平凡的平坦 $U(1)$ 丛 $\mathcal L_i$ 时，如何修改上述公式。此时使用复数系数的 Novikov 环。首先，Floer 复形更自然的写法是
\[
\operatorname{Hom}^i((L_0,\mathcal L_0),(L_1,\mathcal L_1))
=\bigoplus_{\substack{p\in L_0\cap L_1\\I(p)=i}}
\operatorname{Hom}(\mathcal L_{0,p},\mathcal L_{1,p})\otimes\Lambda_{\mathrm{nov}}.
\]
因此，给定 $p\in L_0\cap L_1$ 和 $t\in\operatorname{Hom}(\mathcal L_{0,p},\mathcal L_{1,p})$，便得到一个元素
\[
t_p\in\operatorname{Hom}^*((L_0,\mathcal L_0),(L_1,\mathcal L_1)).
\]

\SourceParagraph{ch08-8.3.3-024}{617}
其次，按照如下方式修改各映射 $m_d$。给定
\[
u\in\mathcal M(p_{0,d};p_{0,1},\ldots,p_{d-1,d};L_0,\ldots,L_d),
\]
对于 $1\leq i\leq d-1$，利用 $\mathcal L_i$ 的联络，沿 $u([x_{i-1,i},x_{i,i+1}])$ 作平行移动，得到对应
\[
H_i^u\colon\mathcal L_{i,p_{i-1,i}}\xrightarrow{\sim}\mathcal L_{i,p_{i,i+1}}.
\]
沿 $u([x_{0,d},x_{0,1}])$ 在 $\mathcal L_0$ 中作平行移动，以及沿 $u([x_{d-1,d},x_{0,d}])$ 在 $\mathcal L_d$ 中作平行移动，则分别得到
\[
H_0^u\colon\mathcal L_{0,p_{0,d}}\xrightarrow{\sim}\mathcal L_{0,p_{0,1}},
\qquad
H_d^u\colon\mathcal L_{d,p_{d-1,d}}\xrightarrow{\sim}\mathcal L_{d,p_{0,d}}.
\]
于是，将 $m_d(t_{p_{d-1,d}},\ldots,t_{p_{0,1}})$ 在点 $p_{0,d}$ 处的系数定义为
\[
\begin{aligned}
&\sum_u(-1)^{s(u)}
\bigl(H_d^u\circ t_{p_{d-1,d}}\circ H_{d-1}^u\circ
t_{p_{d-2,d-1}}\circ\cdots\circ H_0^u\bigr)
\otimes q^{\int u^*\omega}\\
&\hspace{3em}\in\operatorname{Hom}(\mathcal L_{0,p_{0,d}},\mathcal L_{d,p_{0,d}})
\otimes\Lambda_{\mathrm{nov}},
\end{aligned}
\]
其中，求和遍历 $\mathcal M(p_{0,d};p_{0,1},\ldots,p_{d-1,d};L_0,\ldots,L_d)
/\operatorname{PSL}_2(\mathbb R)$ 中的所有 $u$。同一个公式也适用于 $m_1$，以及秩更高的平坦丛。

\SourceParagraph{ch08-8.3.3-025}{617}
到目前为止，这些定义仅适用于 $L_0,\ldots,L_k$ 两两横截的情形，而一般情况下当然未必如此。有两种处理方式。第一种是只定义一个 $A_\infty$ 预范畴；此后若要证明同调镜像对称，就证明它作为预范畴与镜像上的凝聚层导出范畴拟等价。第二种是证明这一构造在 Hamilton 形变下不变。\BookCite{ref421}按如下方式利用这一点：对于每对 Lagrange 子流形 $L_0,L_1$，在 $X$ 上选择一族 Hamilton 函数 $H_t\colon X\to\mathbb R$，$t\in[0,1]$，使相应的时间 $t$ 的 Hamilton 流 $\phi^t$ 满足 $\phi^1(L_0)$ 与 $L_1$ 横截相交。然后仍按照前面的方式定义 $\operatorname{Hom}^*(L_0,L_1)$，但将 $L_0$ 替换为 $\phi^1(L_0)$。这样便能避开横截性问题。特别地，对于单个对象 $L$，也得到了 $\operatorname{Hom}^*(L,L)$。随后可证明，所得 $A_\infty$ 范畴在拟同构意义下与这些选择无关；细节见\BookCite{ref421}。

\SourceParagraph{ch08-8.3.3-026}{618}
虽然前面的一般性讨论较为粗略，但在 $\dim_{\mathbb R}X=2$ 时，我们希望说得更具体一些。对我们而言，这就是 $X$ 为椭圆曲线的情形；下一节将详细处理它。目前，先讨论\BookCite{ref421}、\BookCite{ref002}分析过的符号问题。

\SourceParagraph{ch08-8.3.3-027}{618}
当 $X$ 是椭圆曲线时，任意一维子流形都是 Lagrange 子流形。将 $X$ 写为 $\mathbb C/\Gamma$，其中 $\Gamma$ 为格，便可将 Lagrange 子流形提升到万有覆盖 $\widetilde X=\mathbb C$。因此，给定圆周 $L\subseteq X$ 及其一个提升 $\widetilde L\subseteq\widetilde X$，$\widetilde L$ 或者是圆周（当 $L$ 在 $H_1(X,\mathbb Z)$ 中表示零元时），或者是一条与 $\mathbb R$ 微分同胚的直线。前一种情形没有分次，如\BookRef{fig:8.17}{图17}所示。因此，可以假设 $\widetilde L$ 是一条定向直线。此时，在 $L$ 上取由以圆周 $L$ 为边界的抽象圆盘 $D$ 的唯一自旋结构所诱导的自旋结构。这就是 Neveu--Schwarz 自旋结构，见\BookRef{ex:2.3}{例题2.3}。\BookRef{fig:8.18}{图18}随后给出了交点的指标。

\SourceParagraph{ch08-8.3.3-028}{618}
计算 $A_\infty$ 范畴结构时出现的全纯圆盘，现在可以取为普通的全纯圆盘，也就是 $\mathbb C$ 中的多边形，其边界包含于 $\widetilde L_0\cup\cdots\cup\widetilde L_d$，见\BookRef{fig:8.21}{图21}。图中标出了交点 $p_{0,1},\ldots,p_{3,4}$ 以及 $p_{0,4}$ 的指标。这个图给出一个元素
\[
u\in\mathcal M(p_{0,4};p_{0,1},\ldots,p_{3,4};L_0,\ldots,L_4).
\]
当然，提升 $\widetilde L_i$ 的不同选择会产生不同的圆盘。符号 $(-1)^{s(u)}$ 按如下方式确定：对于每个指标为 $1$ 的点 $p_{i,j}$，$i<j$，有一个符号贡献 $(-1)^{s(p_{i,j})}$；当且仅当圆盘边界的自然定向与 $L_j$ 的给定定向一致时，这个符号为正。例如，在\BookRef{fig:8.21}{图21}中，指标为 $1$ 的交点 $p_{0,1}$ 贡献 $+1$，$p_{2,3}$ 贡献 $-1$，总符号为 $(-1)(+1)=-1$。（在利用带状域，或等价地利用二边形，定义 $\operatorname{Hom}(L_0,L_1)$ 中的 $m_1$ 时，规则是：若 $L_1$ 的给定定向与带状域边界诱导的定向一致，则符号为 $+1$；否则为 $-1$。）

\begin{figure}[H]
\centering
\BookImage{01_Numbered_Figures/Chapter_08/Ch08_Fig21_PDFp618_Bookp608.pdf}
\caption{图中标出了圆盘的标准定向。}\label{fig:8.21}
\end{figure}

\SourceParagraph{ch08-8.3.3-029}{619}
这里还应注意，零维模空间中的圆盘在交点处总有凸角。如\BookRef{fig:8.22}{图22}所示，圆盘的凹角在局部会在一维族中变化：第二幅图中，圆盘边界沿自身折返，而发生折返的点显然可以移动。

\begin{figure}[H]
\centering
\BookImage{01_Numbered_Figures/Chapter_08/Ch08_Fig22_PDFp619_Bookp609.pdf}
\BookNumberOnlyCaption\label{fig:8.22}
\end{figure}

\begin{example*}
\SourceParagraph{ch08-8.3.3-030}{619}
可以在一个简单例子中实施这一过程。取格 $\Gamma=\langle1,\tau\rangle\subseteq\mathbb C$，以及 $L=\mathbb R/\mathbb Z\subseteq X=\mathbb C/\Gamma$。$L$ 的一个提升 $\widetilde L$ 就是 $\mathbb R$。我们希望计算 $\operatorname{Hom}^*(L,L)$。按照上面的讨论，应取 $L_0=L$，而 $L_1$ 为 $L$ 的一个 Hamilton 扰动。在圆柱 $\mathbb C/\mathbb Z$ 中，把 $L$ 的一个小邻域视为圆周 $L$ 的余切丛的零截面邻域，则 $L$ 的一个可能的 Hamilton 形变就是 $L$ 上某个恰当微分的图像，例如
\[
L_1=\{(x,\varepsilon\sin(2\pi x))\mid x\in\mathbb R\},
\]
其中 $\varepsilon>0$ 充分小，使扰动位于这个邻域内。考虑到周期性，有 $L_0\cap L_1=\{p_0,p_1\}$，其中 $p_1=(0,0)$、$p_0=(1/2,0)$；$p_1$ 的指标为 $1$，$p_0$ 的指标为 $0$。\BookRef{fig:8.23}{图23}给出了两个全纯带状域，并标明了它们的符号贡献。因此，$\operatorname{Hom}^*(L_0,L_1)$ 是
\[
[p_0]\Lambda_{\mathrm{nov}}\longrightarrow[p_1]\Lambda_{\mathrm{nov}},
\]
并且
\[
m_1([p_0])=\left(-q^{\int u_1^*\omega}+q^{\int u_2^*\omega}\right)[p_1],
\]
其中，$u_1\colon D\to X$ 和 $u_2\colon D\to X$ 分别对应左边和右边的圆盘。注意到 $\int u_1^*\omega=\int u_2^*\omega$——任何类似的 Hamilton 形变都具有这一性质——可见实际上 $m_1=0$，从而恢复了 $L$ 的以 $\Lambda_{\mathrm{nov}}$ 为系数的上同调 $H^*(L,\Lambda_{\mathrm{nov}})$，正如所预期的那样。
\end{example*}

\begin{figure}[H]
\centering
\BookImage{01_Numbered_Figures/Chapter_08/Ch08_Fig23_PDFp620_Bookp610.pdf}
\BookNumberOnlyCaption\label{fig:8.23}
\end{figure}

\SourceParagraph{ch08-8.3.3-031}{619}
下一节将看到一些更复杂的例子，其中多边形参与椭圆曲线的高阶乘法映射。在此之前，我们最后讨论一下与 $A_\infty$ 范畴有关的代数问题。

\subsection{恒等态射与三角 \texorpdfstring{$A_\infty$}{A∞} 范畴}\label{sec:8.3.4}

\SourceParagraph{ch08-8.3.4-001}{620}
在比较 $X$ 上凝聚层的导出范畴与镜像 $\check X$ 的 Fukaya 范畴时，至少还会遇到两个代数问题，它们似乎使人预期二者不可能同构。第一个问题是，$D^b(X)$ 具有恒等态射，而 $\operatorname{Fuk}(\check X)$ 事先并不具有。第二个问题是，$D^b(X)$ 是三角范畴，而 $\operatorname{Fuk}(\check X)$ 事先并不是。第一个问题结果并不构成障碍：事实上，与最初的预期相反，$\operatorname{Fuk}(\check X)$ 与一个具有恒等态射的 $A_\infty$ 范畴拟同构。第二个问题稍复杂一些。大致而言，与普通范畴相比，三角范畴需要附加数据，即一组满足\BookSectionRef{sec:4.4.3}{4.4.3}节公理的正合三角形。然而，对 $A_\infty$ 范畴而言，成为三角范畴仅仅是一项性质。按照我们目前的定义，尚不知道 $\operatorname{Fuk}(\check X)$ 是否具有这一性质。不过，有一个代数过程可以把它替换为扭曲范畴 $\operatorname{Tw}(\operatorname{Fuk}(\check X))$，后者是包含 $\operatorname{Fuk}(\check X)$ 的三角范畴。介绍这一过程还需要一些工作。

\begin{definition*}
\SourceParagraph{ch08-8.3.4-002}{620}
先讨论恒等态射。设 $\mathcal A$ 是一个（非幺）$A_\infty$ 范畴。恒等态射有几种有用的概念。首先，若非幺的普通范畴 $H^*(\mathcal A)$ 实际上具有恒等态射，即 $H^*(\mathcal A)$ 实际上是一个通常意义的范畴，则称 $\mathcal A$ \emph{上同调幺}，或 \emph{$c$ 幺}。

\SourceParagraph{ch08-8.3.4-003}{620}
另一方面，若对每个对象 $X$ 都存在一个必然唯一的元素 $e_X\in\operatorname{Hom}_{\mathcal A}^0(X,X)$，满足
\[
\begin{gathered}
m_1(e_X)=0,\\
(-1)^{\deg a}m_2(e_{X_1},a)=a=m_2(a,e_{X_0}),
\qquad a\in\operatorname{Hom}_{\mathcal A}(X_0,X_1),\\
\text{任何含有 }e_X\text{ 的高阶乘法映射均为零，}
\end{gathered}
\]
则称 $\mathcal A$ \emph{严格幺}。
\end{definition*}

\SourceParagraph{ch08-8.3.4-004}{620}
首先，$\operatorname{Fuk}(X)$ 实际上是 $c$ 幺的。其原因是：在我们考虑的 Calabi--Yau 情形及本节排除圆盘冒泡的假设下，对于单个分次 Lagrange 子流形 $L$，Floer 上同调 $H^*(\operatorname{Hom}_{\operatorname{Fuk}(X)}(L,L))$ 同构于通常的上同调群 $H^*(L,\Lambda_{\mathrm{nov}})$，而 $1\in H^0(L,\Lambda_{\mathrm{nov}})$ 当然就是恒等元。Piunikhin、Salamon 与 Schwarz 在\BookCite{ref392}中证明了这一事实，虽然是在稍特殊的背景下。随后，Seidel 在\BookCite{ref421}推论 2.14 中证明，任何 $c$ 幺 $A_\infty$ 范畴都与一个严格幺 $A_\infty$ 范畴拟同构。这里不给出细节，但这说明我们不必担心这一问题。因此，从现在起，假设所有 $A_\infty$ 范畴均为严格幺的。

\SourceParagraph{ch08-8.3.4-005}{621}
接下来讨论三角结构。这需要引入若干新的概念。

\BookStatementNumber{14}
\begin{definition}\label{def:8.14}
\SourceParagraph{ch08-8.3.4-006}{621}
设 $\mathcal A$ 是一个 $A_\infty$ 范畴。$\mathcal A$ 上的一个（右）\emph{$A_\infty$ 模} $\mathcal M$，由每个 $X\in\operatorname{Ob}(\mathcal A)$ 对应的分次向量空间 $\mathcal M(X)$，以及对每个 $d\geq1$ 给定的下列乘法映射组成：
\setcounter{equation}{8}
\begin{equation}\label{eq:8.9}
\begin{aligned}
m_d^{\mathcal M}\colon{}&\mathcal M(X_{d-1})\otimes
\operatorname{Hom}_{\mathcal A}(X_{d-2},X_{d-1})\otimes\cdots
\otimes\operatorname{Hom}_{\mathcal A}(X_0,X_1)\\
&\longrightarrow\mathcal M(X_0)[2-d].
\end{aligned}
\end{equation}
这些映射应满足关系
\[
\begin{aligned}
&\sum_n(-1)^{\deg a_1+\cdots+\deg a_n-n}
m_{n+1}^{\mathcal M}\bigl(\\
&\qquad m_{d-n}^{\mathcal M}(b,a_{d-1},\ldots,a_{n+1}),a_n,\ldots,a_1\bigr)\\
&\quad+\sum_{m,n}(-1)^{\deg a_1+\cdots+\deg a_n-n}\,
m_{d-m+1}^{\mathcal M}\bigl(b,a_{d-1},\ldots,a_{n+m+1},\\
&\qquad\quad m_m^{\mathcal A}(a_{n+m},\ldots,a_{n+1}),a_n,\ldots,a_1\bigr)=0.
\end{aligned}
\]

\SourceParagraph{ch08-8.3.4-007}{621}
第二个和遍历 $n+m<d$ 的指标。例如，当 $d=1$ 时，这个关系成为 $(m_1^{\mathcal M})^2=0$，即 $m_1^{\mathcal M}$ 使 $\mathcal M(X)$ 成为链复形。当 $d=2$ 时，它为
\[
m_1^{\mathcal M}(m_2^{\mathcal M}(b,a_1))
+(-1)^{\deg a_1-1}m_2^{\mathcal M}(m_1^{\mathcal M}(b),a_1)
+m_2^{\mathcal M}(b,m_1^{\mathcal A}(a_1))=0,
\]
而当 $d=3$ 时，它为
\[
\begin{aligned}
&m_1^{\mathcal M}(m_3^{\mathcal M}(b,a_2,a_1))
+(-1)^{\deg a_1-1}m_2^{\mathcal M}(m_2^{\mathcal M}(b,a_2),a_1)\\
&\quad+(-1)^{\deg a_1+\deg a_2-2}
m_3^{\mathcal M}(m_1^{\mathcal M}(b),a_2,a_1)\\
&\quad+m_3^{\mathcal M}(b,a_2,m_1^{\mathcal A}(a_1))\\
&\quad+(-1)^{\deg a_1-1}m_3^{\mathcal M}(b,m_1^{\mathcal A}(a_2),a_1)
+m_2^{\mathcal M}(b,m_2^{\mathcal A}(a_2,a_1))=0.
\end{aligned}
\]
与原来的 $A_\infty$ 关系一样，后面这两个展开式中的第一个说明 $m_1^{\mathcal M}$ 对模乘法 $m_2^{\mathcal M}$ 满足 Leibniz 法则，而 $m_3^{\mathcal M}$ 则衡量这一模乘法偏离结合律的程度。
\end{definition}

\SourceParagraph{ch08-8.3.4-008}{621}
$\mathcal A$ 上的 $A_\infty$ 模本身也构成一个 $A_\infty$ 范畴 $\operatorname{Mod}(\mathcal A)$。需要定义 $A_\infty$ 模之间的态射概念，我们称之为\emph{预同态}。若 $\mathcal M_1,\mathcal M_2$ 是两个 $A_\infty$ 模，则次数为 $\deg t$ 的预同态 $t\colon\mathcal M_1\to\mathcal M_2$，由对每个 $d\geq1$ 给定的一列映射组成：
\setcounter{equation}{9}
\begin{equation}\label{eq:8.10}
\begin{aligned}
t_d\colon{}&\mathcal M_1(X_{d-1})\otimes
\operatorname{Hom}_{\mathcal A}(X_{d-2},X_{d-1})\otimes\cdots
\otimes\operatorname{Hom}_{\mathcal A}(X_0,X_1)\\
&\longrightarrow\mathcal M_2(X_0)[\deg t-d+1].
\end{aligned}
\end{equation}

\SourceParagraph{ch08-8.3.4-009}{622}
随后，在预同态的分次向量空间 $\operatorname{Hom}_{\operatorname{Mod}(\mathcal A)}(\mathcal M_1,\mathcal M_2)$ 上，按下式定义 $m_1^{\operatorname{Mod}(\mathcal A)}$：
\[
\begin{aligned}
&\bigl(m_1^{\operatorname{Mod}(\mathcal A)}(t)_d\bigr)(b,a_{d-1},\ldots,a_1)\\
&\quad=\sum_n(-1)^*m_{n+1}^{\mathcal M_2}
\bigl(t_{d-n}(b,a_{d-1},\ldots,a_{n+1}),a_n,\ldots,a_1\bigr)\\
&\qquad+\sum_n(-1)^*t_{n+1}
\bigl(m_{d-n}^{\mathcal M_1}(b,a_{d-1},\ldots,a_{n+1}),a_n,\ldots,a_1\bigr)\\
&\qquad+\sum_{m,n}(-1)^*t_{d-m+1}
\bigl(b,a_{d-1},\ldots,a_{n+m+1},
m_m^{\mathcal A}(a_{n+m},\ldots,a_{n+1}),a_n,\ldots,a_1\bigr).
\end{aligned}
\]
这里，$*$ 表示表达式
\[
\deg a_{n+1}+\cdots+\deg a_{d-1}+\deg b-d+n+1.
\]
此外，还需要定义复合：给定 $t^1\colon\mathcal M_0\to\mathcal M_1$ 和 $t^2\colon\mathcal M_1\to\mathcal M_2$，令
\[
\begin{aligned}
&\bigl(m_2^{\operatorname{Mod}(\mathcal A)}(t^2,t^1)\bigr)_d
(b,a_{d-1},\ldots,a_1)\\
&\quad=\sum_n(-1)^*t^2_{n+1}
\bigl(t^1_{d-n}(b,a_{d-1},\ldots,a_{n+1}),a_n,\ldots,a_1\bigr).
\end{aligned}
\]

\SourceParagraph{ch08-8.3.4-010}{622}
读者可以在低次情形将这些式子展开，以再次体会它们的含义。幸运的是，此时所有高阶乘法映射 $m_d^{\operatorname{Mod}(\mathcal A)}$（$d>2$）均为零，所以得到的是一个微分分次范畴。若 $m_1^{\operatorname{Mod}(\mathcal A)}(t)=0$，则称 $t$ 为\emph{同态}。这个范畴是严格幺的，其中 $\operatorname{id}\in\operatorname{Hom}_{\operatorname{Mod}(\mathcal A)}(\mathcal M,\mathcal M)$ 由
\[
\operatorname{id}_1(b)=(-1)^{\deg b}b,\qquad
\operatorname{id}_d=0\quad(d>1)
\]
给出。所谓\emph{同构}，是指在 $H^*(\operatorname{Mod}(\mathcal A))$ 中为同构的同态。

\SourceParagraph{ch08-8.3.4-011}{622}
现在可以定义称为 \emph{Yoneda 嵌入}的 $A_\infty$ 函子 $\operatorname{Yon}\colon\mathcal A\to\operatorname{Mod}(\mathcal A)$。它的定义很简单：对于 $Y\in\operatorname{Ob}(\mathcal A)$，令 $\operatorname{Yon}(Y)$ 为模
\[
X\longmapsto\operatorname{Hom}_{\mathcal A}(X,Y).
\]
当然，后者已经是由 $m_1^{\mathcal A}$ 给出的链复形。为了定义模结构，需要给出\BookEqRef{eq:8.9}{8.9}中的映射。但在此情况下，$\mathcal M(X_{d-1})=\operatorname{Hom}_{\mathcal A}(X_{d-1},Y)$，且 $\mathcal M(X_0)=\operatorname{Hom}_{\mathcal A}(X_0,Y)$，所以这些映射直接由 $m_d^{\mathcal A}$ 给出。这样便定义了函子 $\operatorname{Yon}$ 的对象映射，将 $\mathcal A$ 的对象映为 $\operatorname{Mod}(\mathcal A)$ 的对象。不过，还必须定义高阶映射
\[
\begin{aligned}
\operatorname{Yon}_d\colon{}&\operatorname{Hom}_{\mathcal A}(X_{d-1},X_d)
\otimes\cdots\otimes\operatorname{Hom}_{\mathcal A}(X_0,X_1)\\
&\longrightarrow
\operatorname{Hom}_{\operatorname{Mod}(\mathcal A)}
(\operatorname{Yon}(X_0),\operatorname{Yon}(X_d))[1-d].
\end{aligned}
\]

\SourceParagraph{ch08-8.3.4-012}{622}
因此，给定 $a_1,\ldots,a_d$，需要给出一个预同态，即一组映射
\[
\begin{aligned}
t_n\colon&\operatorname{Yon}(X_0)(Y_{n-1})\otimes
\operatorname{Hom}_{\mathcal A}(Y_{n-2},Y_{n-1})\otimes\cdots
\otimes\operatorname{Hom}_{\mathcal A}(Y_0,Y_1)\\
&\longrightarrow\operatorname{Yon}(X_d)(Y_0)
\left[2-d-n+\sum_i\deg a_i\right],
\end{aligned}
\]
也就是
\[
\begin{aligned}
(\operatorname{Yon}_d)_n\colon&\operatorname{Hom}_{\mathcal A}(Y_{n-1},X_0)
\otimes\operatorname{Hom}_{\mathcal A}(Y_{n-2},Y_{n-1})\otimes\cdots
\otimes\operatorname{Hom}_{\mathcal A}(Y_0,Y_1)\\
&\longrightarrow\operatorname{Hom}_{\mathcal A}(Y_0,X_d)
\left[2-d-n+\sum_i\deg a_i\right].
\end{aligned}
\]
有一个显然的选择：
\[
\begin{aligned}
&\bigl(\operatorname{Yon}_d(a_d,\ldots,a_1)\bigr)_n(b,b_{n-1},\ldots,b_1)\\
&\hspace{3em}=m_{d+n}^{\mathcal A}(a_d,\ldots,a_1,b,b_{n-1},\ldots,b_1).
\end{aligned}
\]
这定义了一个 $A_\infty$ 函子 $\operatorname{Yon}$，读者可以自行验证。

\SourceParagraph{ch08-8.3.4-013}{623}
类似于普通范畴论中的情形，若存在同构 $t\colon\operatorname{Yon}(Y)\to\mathcal M$，就称对象 $Y\in\operatorname{Ob}(\mathcal A)$ \emph{拟表示} $A_\infty$ 模 $\mathcal M$。这种处理的优点在于，$\operatorname{Mod}(\mathcal A)$ 中的代数操作容易定义；当这些操作的结果可拟表示时，它们便转化为 $\mathcal A$ 中的概念。有两个例子对我们至关重要。

\SourceParagraph{ch08-8.3.4-014}{623}
首先是\emph{平移函子}：给定 $A_\infty$ 模 $\mathcal M$，定义 $S\mathcal M$ 为
\[
(S\mathcal M)(Y)=\mathcal M(Y)[1].
\]
若对所有 $Y\in\operatorname{Ob}(\mathcal A)$，$S(\operatorname{Yon}(Y))$ 均可由某个 $SY\in\operatorname{Ob}(\mathcal A)$ 拟表示，便得到平移函子 $S\colon\mathcal A\to\mathcal A$。例如，在 Fukaya 范畴中，只需将分次加 $1$，就能实现平移函子。

\SourceParagraph{ch08-8.3.4-015}{623}
对我们更重要的是\emph{映射锥}构造。设 $Y_0,Y_1\in\operatorname{Ob}(\mathcal A)$，$c\in\operatorname{Hom}_{\mathcal A}^0(Y_0,Y_1)$，满足 $m_1^{\mathcal A}(c)=0$。可以将 $c$ 的抽象映射锥定义为如下 $A_\infty$ 模 $\operatorname{Cone}(c)$：
\[
\operatorname{Cone}(c)(X):=
\operatorname{Hom}_{\mathcal A}(X,Y_0)[1]\oplus
\operatorname{Hom}_{\mathcal A}(X,Y_1),
\]
\[
\begin{aligned}
&m_d^{\operatorname{Cone}(c)}((b_0,b_1),a_{d-1},\ldots,a_1)\\
&\quad=\bigl(m_d^{\mathcal A}(b_0,a_{d-1},\ldots,a_1),\,
m_d^{\mathcal A}(b_1,a_{d-1},\ldots,a_1)\\
&\hspace{6em}
+m_{d+1}^{\mathcal A}(c,b_0,a_{d-1},\ldots,a_1)\bigr).
\end{aligned}
\]
若 $\mathcal A$ 中存在一个拟表示 $\operatorname{Cone}(c)$ 的对象，也将它记为 $\operatorname{Cone}(c)$。

\SourceParagraph{ch08-8.3.4-016}{623}
模 $\operatorname{Cone}(c)$ 带有典范的预同态
\[
\begin{aligned}
\iota&\in\operatorname{Hom}_{\operatorname{Mod}(\mathcal A)}^0
(\operatorname{Yon}(Y_1),\operatorname{Cone}(c)),\\
\pi&\in\operatorname{Hom}_{\operatorname{Mod}(\mathcal A)}^1
(\operatorname{Cone}(c),\operatorname{Yon}(Y_0)),
\end{aligned}
\]
由
\[
\iota_1(b_1)=(0,(-1)^{\deg b_1}b_1),\qquad
\pi_1(b_0,b_1)=(-1)^{\deg b_0-1}b_0
\]
给出，而这些预同态的所有高阶部分均为零。由此在 $H^*(\operatorname{Mod}(\mathcal A))$ 中得到一个“三角形”：
\setcounter{equation}{10}
\begin{figure}[H]
\begin{equation}\label{eq:8.11}
\BookImage{02_Unnumbered_Diagrams/Chapter_08/Ch08_Diagram_PDFp623_01_Bookp613.pdf}
\end{equation}
\end{figure}

\SourceParagraph{ch08-8.3.4-017}{624}
这里，$[\cdot]$ 表示同态在 $H^*(\operatorname{Mod}(\mathcal A))$ 中的上同调类。由于 $m_1^{\operatorname{Mod}(\mathcal A)}$ 作用于 $\operatorname{Yon}_1(c)$、$\iota$ 或 $\pi$ 时均为零，这些上同调类有意义，读者可以自行验证。请记住，$[\pi]$ 的次数为一。

\SourceParagraph{ch08-8.3.4-018}{624}
现在可以定义 $\mathcal A$ 中\emph{正合三角形}的概念。它是 $H^*(\mathcal A)$ 中如下形式的图表：
\begin{figure}[H]
\centering
\BookImage{02_Unnumbered_Diagrams/Chapter_08/Ch08_Diagram_PDFp624_01_Bookp614.pdf}
\end{figure}
它在 Yoneda 嵌入下变为与\BookEqRef{eq:8.11}{8.11}同构的图表，其中 $c_1=c$。

\BookStatementNumber{15}
\begin{definition}\label{def:8.15}
\SourceParagraph{ch08-8.3.4-019}{624}
若 $A_\infty$ 范畴 $\mathcal A$ 满足以下条件，则称它为\emph{三角} $A_\infty$ 范畴：
\begin{enumerate}
\item $H^0(\mathcal A)$ 中的每个态射 $[c_1]$ 都可以扩充为一个正合三角形。
\item 存在平移函子 $S\colon\mathcal A\to\mathcal A$。
\item 对每个 $Y\in\operatorname{Ob}(\mathcal A)$，均存在对象 $\widetilde Y\in\operatorname{Ob}(\mathcal A)$，使 $S\widetilde Y\cong Y$ 在 $H^0(\mathcal A)$ 中成立。
\end{enumerate}
\end{definition}

\SourceParagraph{ch08-8.3.4-020}{624}
三角 $A_\infty$ 范畴的一个例子是 $D^b_\infty(X)$：对于它的对象 $M$，平移函子就是\BookSectionRef{sec:4.4.3}{4.4.3}节定义的通常的复形平移 $M\mapsto M[1]$。此外，这个范畴中的映射锥也按照\BookSectionRef{sec:4.4.3}{4.4.3}节的方式定义。

\SourceParagraph{ch08-8.3.4-021}{624}
另一方面，尚不知道 Fukaya 范畴是否是三角的。至少预计需要把浸入而非嵌入的 Lagrange 子流形也加入 Fukaya 范畴，而这会给 Floer 同调的定义带来技术困难。所以，仍需解决这一问题。解决办法是\emph{扭曲构造}：给定 $A_\infty$ 范畴 $\mathcal A$，必要时先加入对象的形式平移，再定义新的 $A_\infty$ 范畴 $\operatorname{Tw}(\mathcal A)$，即\emph{扭曲复形范畴}；它是三角的。

\BookStatementNumber{16}
\begin{definition}\label{def:8.16}
\SourceParagraph{ch08-8.3.4-022}{624}
设 $\mathcal A$ 是一个 $A_\infty$ 范畴。一个\emph{扭曲复形}，由 $\mathcal A$ 的一列对象 $X_1,\ldots,X_n$（在需要时包括它们的形式平移），以及一个由一次态射组成的严格下三角矩阵 $\Delta=(\delta_{i,j})_{1\leq i,j\leq n}$ 构成，其中 $\delta_{i,j}\in\operatorname{Hom}^1(X_j,X_i)$，并满足
\[
\sum_{k\geq1}m_k(\Delta,\ldots,\Delta)=0.
\]
这里按通常的矩阵乘法规则，理解 $m_k$ 作用于态射矩阵的含义。因而，利用严格下三角条件，上式等价于对每对 $p>q$ 都有
\[
\sum_{\substack{k\geq1\\q=i_0<i_1<\cdots<i_k=p}}
m_k(\delta_{i_k,i_{k-1}},\ldots,\delta_{i_1,i_0})=0.
\]

\SourceParagraph{ch08-8.3.4-023}{625}
注意，这个看似无穷的和因此只有有限个非零项。记 $X=(X_1,\ldots,X_n)$，$\Delta=\Delta^X$，于是将扭曲对象记为 $(X,\Delta^X)$。
\end{definition}

\SourceParagraph{ch08-8.3.4-024}{625}
$\mathcal A$ 的扭曲范畴是以 $\mathcal A$ 的扭曲复形为对象的 $A_\infty$ 范畴，其态射空间为
\[
\operatorname{Hom}_{\operatorname{Tw}(\mathcal A)}((X,\Delta^X),(Y,\Delta^Y))
=\bigoplus_{i,j}\operatorname{Hom}_{\mathcal A}(X_i,Y_j).
\]
若 $(X^0,\Delta^{X^0}),\ldots,(X^d,\Delta^{X^d})$ 是一列扭曲复形，并且
\[
f^i:=(f_{j,k}^i)\in\operatorname{Hom}_{\operatorname{Tw}(\mathcal A)}
((X^{i-1},\Delta^{X^{i-1}}),(X^i,\Delta^{X^i})),
\]
则映射 $m_d^{\operatorname{Tw}(\mathcal A)}$ 为
\[
\begin{aligned}
&m_d^{\operatorname{Tw}(\mathcal A)}(f^d,\ldots,f^1)\\
&\quad=\sum_{j_0,\ldots,j_d}
m_{d+j_0+\cdots+j_d}^{\mathcal A}\bigl(
\underbrace{\Delta^{X^d},\ldots,\Delta^{X^d}}_{j_d},f^d,\\
&\hspace{6em}
\underbrace{\Delta^{X^{d-1}},\ldots,\Delta^{X^{d-1}}}_{j_{d-1}},
f^{d-1},\ldots,f^1,
\underbrace{\Delta^{X^0},\ldots,\Delta^{X^0}}_{j_0}\bigr).
\end{aligned}
\]
需要作一点说明：首先，求和遍历所有 $j_0,\ldots,j_d\geq0$；其次，$\Delta^{X^i}$ 在表达式中出现 $j_i$ 次；第三，与前面一样，应将该表达式理解为矩阵乘法。我们将验证这些 $m_d$ 满足 $A_\infty$ 关系的工作留给读者！

\SourceParagraph{ch08-8.3.4-025}{625}
最简单的例子是直和，此时 $\Delta^X=0$，所以对象 $((X_1,\ldots,X_n),0)$ 表示 $X_1\oplus\cdots\oplus X_n$。一个稍不平凡的扭曲复形例子是映射锥构造。假设 $A_\infty$ 范畴 $\mathcal A$ 具有平移函子 $S$，Fukaya 范畴就是如此。给定满足 $m_1^{\mathcal A}(c)=0$ 的 $c\in\operatorname{Hom}_{\mathcal A}^0(Y_0,Y_1)$，得到扭曲复形
\[
C=\left((SY_0,Y_1),\begin{pmatrix}0&0\\-S(c)&0\end{pmatrix}\right),
\]
其中，利用平移函子 $S$，把 $c$ 视为定义了 $\operatorname{Hom}_{\mathcal A}^1(SY_0,Y_1)$ 中的一个元素 $S(c)$。注意，存在典范态射
\[
\operatorname{Hom}^0(Y_1,C)=
\operatorname{Hom}^0(Y_1,SY_0)\oplus\operatorname{Hom}^0(Y_1,Y_1)
\]
中的 $(0,\operatorname{id}_{Y_1})$，以及
\[
\begin{aligned}
\operatorname{Hom}^1(C,Y_0)
&=\operatorname{Hom}^1(SY_0,Y_0)\oplus\operatorname{Hom}^1(Y_1,Y_0)\\
&=\operatorname{Hom}^0(Y_0,Y_0)\oplus\operatorname{Hom}^1(Y_1,Y_0)
\end{aligned}
\]
中的 $(\operatorname{id}_{Y_0},0)$。读者同样可以验证，这些映射在 $\operatorname{Tw}(\mathcal A)$ 中给出一个正合三角形，即 $C$ 在 $\operatorname{Tw}(\mathcal A)$ 中表示 $c$ 的映射锥。

\SourceParagraph{ch08-8.3.4-026}{625}
事实上，稍作推广就可看出，$\operatorname{Tw}(\mathcal A)$ 本身是三角的。给定闭态射
\[
c\in\operatorname{Hom}_{\operatorname{Tw}(\mathcal A)}^0
((X^0,\Delta^{X^0}),(X^1,\Delta^{X^1})),
\]
将 $\operatorname{Cone}(c)$ 构造为
\[
C=\left((SX^0,X^1),
\begin{pmatrix}\Delta^{SX^0}&0\\-S(c)&\Delta^{X^1}\end{pmatrix}\right),
\]
其中左上块为将 $X^0$ 连同其扭曲数据一并平移所得的矩阵。

\SourceParagraph{ch08-8.3.4-027}{626}
可以将 $\operatorname{Tw}(\mathcal A)$ 视为 $\mathcal A$ 的导出范畴，并把它写为 $D^b(\mathcal A)$。

\SourceParagraph{ch08-8.3.4-028}{626}
遗憾的是，这还不完全是我们所需要的范畴。$X$ 上凝聚层的导出范畴 $D^b(X)$ 还有一个 $\operatorname{Fuk}(\check X)$ 或 $D^b(\operatorname{Fuk}(\check X))$ 未必具有的性质：若 $F$ 是 $D^b(X)$ 的对象，$p\colon F\to F$ 是满足 $p^2=p$ 的态射，则可以把 $p$ 的像看作 $D^b(X)$ 中的另一个对象。这样的态射 $p$ 称为\emph{幂等态射}，而 $D^b(X)$ 是\emph{分裂闭}的；实质上，这意味着 $D^b(X)$ 包含所有幂等态射的像。

\SourceParagraph{ch08-8.3.4-029}{626}
一般而言，对于给定的线性范畴，可以构造一个包含所有幂等态射的像的更大范畴，即 \emph{Karoubi 完备化}。这可以推广到 $A_\infty$ 范畴，从 $A_\infty$ 范畴 $\mathcal A$ 得到一个分裂闭范畴。我们略去构造细节；感兴趣的读者可参阅\BookCite{ref421}第 4 节。该构造与前面看到的构造具有相似的特点。特别地，以 $D^\pi(\mathcal A)$ 表示将这一过程施加于 $D^b(\mathcal A)$ 所得的结果。可以证明它是一个分裂闭的三角 $A_\infty$ 范畴。完成这一步后，我们所希望的同调镜像对称的严格表述就成为：

\BookStatementNumber{17}
\begin{conjecture}\label{conj:8.17}
\SourceParagraph{ch08-8.3.4-030}{626}
设 $X$ 和 $\check X$ 为一对互为镜像的 Calabi--Yau 流形。则 $D^b_\infty(X)$ 与 $D^\pi(\operatorname{Fuk}(\check X))$ 作为 $A_\infty$ 范畴拟等价。
\end{conjecture}
